\documentclass[12pt]{article}
\usepackage[T1]{fontenc}
\usepackage[utf8]{inputenc}
\usepackage{lmodern}
\usepackage{textcomp}
\usepackage{enumitem}
\usepackage{geometry}
\usepackage{graphicx}
\usepackage{amsmath, amssymb}
\geometry{margin=1in}
\begin{document}
\begin{center}
Sociology - Graduate Level Assessment 12 \
Total Marks: 40 \
Answer all questions.
\end{center}

\section*{Multiple Choice (10 marks)}
\begin{enumerate}[label=\arabic*.]
\item In idealized views of science, the experimental method is said to involve:
\begin{enumerate}[label=(\alph*)]
    \item testing out new research methods to see which one works best
    \item isolating and measuring the effect of one variable upon another
    \item using personal beliefs and values to decide what to study
    \item interpreting data subjectively, drawing on theoretical paradigms
\end{enumerate}
\item Goldthorpe identified the 'service class' as:
\begin{enumerate}[label=(\alph*)]
    \item those in non-manual occupations, exercising authority on behalf of the state
    \item people working in consultancy firms who were recruited by big businesses
    \item the young men and women employed in domestic service in the nineteenth century
    \item those who had worked in the armed services
\end{enumerate}
\item Which of the following is not identified as a new form of community?
\begin{enumerate}[label=(\alph*)]
    \item ethnic communities, based on shared identity and experiences of discrimination
    \item gay villages, which are formed in certain parts of large cities
    \item sociological communities, formed by unpopular lecturers
    \item virtual communities that exist only in cyberspace
\end{enumerate}
\item Sociology can be considered a social science because:
\begin{enumerate}[label=(\alph*)]
    \item its theories are logical, explicit and supported by empirical evidence
    \item sociologists collect data in a relatively objective and systematic way
    \item ideas and research findings are scrutinized by other sociologists
    \item all of the above
\end{enumerate}
\item The introduction of market principles to educational policy in the 1980 s resulted in:
\begin{enumerate}[label=(\alph*)]
    \item more funding for students in higher education
    \item the delegation of power and budgetary control to LEAs
    \item a reduction in parental choice of school
    \item increased state regulation through national testing and inspections
\end{enumerate}
\end{enumerate}

\section*{True / False (10 marks)}
\textit{Answer True or False}
\begin{enumerate}[label=\arabic*.]
\setcounter{enumi}{5}
\item True or False: Dahrendorf, Rex, and Habermas concentrated on concepts of power, domination, and conflict within social structures.
\item True or False: Knowledge is purely objective and does not reflect the personal values of theorists or sociologists.
\item True or False: Cultural restructuring primarily focuses on preserving traditional landscapes without any economic transformation or media involvement.
\item True or False: Lombroso believed that all criminal behavior was solely a result of socialization into a criminal underworld.
\item True or False: Women may seek cosmetic surgery due to perceiving certain body parts as stigmatizing in relation to cultural ideals of beauty.
\end{enumerate}

\section*{Long Form Response (20 marks)}
\textit{Answer each question with a clear and complete explanation.}
\begin{enumerate}[label=\arabic*.]
\setcounter{enumi}{10}
\item Discuss Marx's assertion that religion would vanish following a socialist revolution. Analyze the relationship between capitalism, ideology, and the role of religion in society.
\item Discuss the characteristics of totalitarian societies, highlighting the implications of restricted freedoms, particularly the absence of freedom of movement for citizens.
\end{enumerate}

\vfill
\noindent\textit{End of Paper}
\end{document}